\documentclass[11pt,a4paper]{article}

\usepackage[T1]{fontenc}
\usepackage[utf8]{inputenc}
\usepackage{amsmath,amssymb,amsthm,mathtools,bm}
\usepackage{newtxtext,newtxmath}
\usepackage[margin=27mm]{geometry}
\usepackage{microtype}
\usepackage{booktabs,longtable,array,tabularx}
\usepackage{enumitem}
\usepackage{xcolor}
\usepackage[most]{tcolorbox}
\usepackage{listings}
\usepackage{url}
\usepackage{hyperref}
\usepackage[nameinlink,capitalise,noabbrev]{cleveref}

\definecolor{DeepBlue}{RGB}{23,60,104}
\definecolor{SoftBlue}{RGB}{238,245,252}
\definecolor{DeepGreen}{RGB}{38,92,62}
\definecolor{SoftGreen}{RGB}{240,249,243}
\definecolor{DeepRed}{RGB}{130,42,42}
\definecolor{SoftRed}{RGB}{252,242,242}
\definecolor{CodeGray}{RGB}{247,247,247}

\hypersetup{
  colorlinks=true,
  linkcolor=DeepBlue,
  citecolor=DeepGreen,
  urlcolor=DeepBlue,
  pdftitle={Inverting the K-Function at Infinity},
  pdfauthor={Research report for the ProveIt project}
}

\setlist[itemize]{topsep=3pt,itemsep=2pt,parsep=0pt}
\setlist[enumerate]{topsep=3pt,itemsep=2pt,parsep=0pt}
\allowdisplaybreaks
\numberwithin{equation}{section}

\newtheorem{theorem}{Theorem}[section]
\newtheorem{proposition}[theorem]{Proposition}
\newtheorem{lemma}[theorem]{Lemma}
\newtheorem{corollary}[theorem]{Corollary}
\newtheorem{definition}[theorem]{Definition}
\newtheorem{algorithm}[theorem]{Algorithm}
\newtheorem{conjecture}[theorem]{Conjecture}
\theoremstyle{remark}
\newtheorem{remark}[theorem]{Remark}
\newtheorem{example}[theorem]{Example}

\newtcolorbox{mainresult}{
  enhanced,
  breakable,
  colback=SoftBlue,
  colframe=DeepBlue,
  boxrule=0.8pt,
  arc=1.5mm,
  title=Main result,
  fonttitle=\bfseries
}
\newtcolorbox{practicalbox}{
  enhanced,
  breakable,
  colback=SoftGreen,
  colframe=DeepGreen,
  boxrule=0.8pt,
  arc=1.5mm,
  title=Practical formula,
  fonttitle=\bfseries
}
\newtcolorbox{warningbox}{
  enhanced,
  breakable,
  colback=SoftRed,
  colframe=DeepRed,
  boxrule=0.8pt,
  arc=1.5mm,
  title=Important ordering point,
  fonttitle=\bfseries
}

\lstset{
  basicstyle=\ttfamily\small,
  backgroundcolor=\color{CodeGray},
  frame=single,
  rulecolor=\color{black!20},
  breaklines=true,
  columns=fullflexible,
  keepspaces=true,
  showstringspaces=false,
  upquote=true,
  numbers=left,
  numberstyle=\tiny\color{black!55},
  xleftmargin=1.8em,
  framexleftmargin=1.4em
}

\DeclareMathOperator{\Log}{Log}
\DeclareMathOperator{\ord}{ord}
\DeclareMathOperator{\Res}{Res}
\newcommand{\e}{\mathrm e}
\newcommand{\ii}{\mathrm i}
\newcommand{\Aconst}{A}
\newcommand{\Kinv}{K_{+}^{-1}}
\newcommand{\asympseries}{\mathrel{\sim}}
\newcommand{\bigO}{\mathcal O}
\newcommand{\smallo}{o}
\newcommand{\dd}{\,\mathrm d}
\newcommand{\coeff}[2]{[{}#1^{#2}]}

\title{\textbf{Inverting the K-Function at Infinity}\\[4pt]
\large Lambert--W Normalization, All-Orders Transseries,\\
and a General Theory of Power--Logarithmic Reversion}
\author{Research report prepared for the ProveIt project}
\date{3 September 2026}

\begin{document}
\maketitle

\begin{abstract}
The K-function is the complex interpolation of the hyperfactorial,
\(K(n+1)=\prod_{j=1}^{n}j^j\).  Its logarithm grows like
\(\tfrac12x^2\log x-\tfrac14x^2\), so its large positive inverse is not
naturally expanded in ordinary powers of \(1/\log Z\).  The correct first
step is to invert the dominant power--logarithmic phase exactly by the
Lambert W-function.  After the unique centering \(r=x-\tfrac12\), we prove
that, with \(Y=\log Z\),
\[
 w=W_0(4Y/\e),\qquad
 \rho=2\sqrt{Y/w}=\exp\!\left(\frac{w+1}{2}\right),\qquad
 L=\log\rho,
\]
the principal inverse has the all-orders transseries
\[
 K_{+}^{-1}(Z)\sim \frac12+\rho+
 \sum_{n\ge1}V_n(L)\rho^{1-2n}.
\]
The coefficients \(V_n\) are rational Laurent polynomials in \(L\), with
polynomial dependence on \(C=\log A-\tfrac18\), and are generated by an
explicit triangular recurrence.  We compute the first coefficients and
prove the asymptotic validity of every finite truncation.

A second, substantially faster normal form is obtained from
\(T=x^2-x+\tfrac16\).  In this coordinate,
\[
 \log K(x)\sim \frac14T(\log T-1)+C+
 \sum_{n\ge1}d_nT^{-n},
\]
so the leading inverse is
\[
 K_{+}^{-1}(Z)\approx \frac12+
 \sqrt{\frac1{12}+
 \frac{4(Y-C)}{W_0(4(Y-C)/\e)}}.
\]
An all-orders correction series for \(T\) is derived.  Numerical tests show
that this accelerated formula and its first corrections are highly accurate
even at moderate arguments.

The K-function is then used as a model for a general inversion theorem for
transseries with dominant part \(a x^p(\log x-b)\).  We give the exact
Lambert--W normalizer, a complete formal recurrence over a slow coefficient
field, residual-to-error estimates, Newton acceleration, branch formulas in
the complex plane, and the transport rule for exponentially small sectors.
The derivations are self-contained at the asymptotic level; no claim of
literature priority is made, and the results have not been machine-checked.
\end{abstract}

\tableofcontents
\newpage

\begin{mainresult}
Let \(K_+\) be the restriction of the K-function to \([2,\infty)\), and let
\(Z\to+\infty\).  Put
\[
Y=\log Z,\qquad C=\log A-\frac18,
\qquad w=W_0\!\left(\frac{4Y}{\e}\right),
\qquad \rho=2\sqrt{\frac{Y}{w}},
\qquad L=\frac{w+1}{2}.
\]
Then
\[
\boxed{
\Kinv(Z)\sim \frac12+\rho+
\frac{V_1(L)}{\rho}+
\frac{V_2(L)}{\rho^3}+
\frac{V_3(L)}{\rho^5}+\cdots
}
\]
where
\[
E=C-\frac{L}{24},\qquad
V_1=-\frac{E}{L}=\frac1{24}-\frac{C}{L},
\]
\[
V_2=-\frac1L\left(\frac{L+1}{2}V_1^2-\frac{V_1}{24}-\frac7{5760}\right),
\]
and
\[
V_3=-\frac1L\left((L+1)V_1V_2+\frac{V_1^3}{6}
-\frac1{24}\left(V_2-\frac{V_1^2}{2}\right)
+\frac7{2880}V_1+\frac{31}{161280}\right).
\]
All subsequent coefficients are determined uniquely by the recurrence in
\cref{thm:canonical-recurrence}.

For computation, the accelerated coordinate
\[
T=x^2-x+\frac16
\]
is preferable.  With
\[
\eta=Y-C,
\qquad \omega=W_0\!\left(\frac{4\eta}{\e}\right),
\qquad \tau=\frac{4\eta}{\omega},
\qquad S=\log\tau=\omega+1,
\]
one has
\[
\boxed{
\Kinv(Z)\sim \frac12+
\sqrt{\frac1{12}+\tau\left(1+
\sum_{n\ge2}U_n(S)\tau^{-n}\right)}
}
\]
with
\[
U_2=\frac{1}{120S},\qquad
U_3=-\frac{31}{22680S},
\]
\[
U_4=\frac{1030S^2-126S-63}{1814400S^3},\qquad
U_5=-\frac{16110S^2-1023S-341}{29937600S^3}.
\]
\end{mainresult}

\section{Introduction and orientation}

\subsection{Why this inverse requires a transseries}

The hyperfactorial
\[
H(n)=\prod_{j=1}^{n}j^j
\]
grows faster than an ordinary exponential but slower than a function such as
\(\exp(n^3)\).  Its continuous interpolation by the K-function satisfies
\[
\log K(x)=\frac12x^2\log x-\frac14x^2+\text{lower-order terms}.
\]
Consequently, solving \(K(x)=Z\) begins with the nonlinear balance
\[
\frac12x^2\log x\asymp \log Z.
\]
This is already outside the scope of a pure power series in \(1/\log Z\).
The leading scale contains a Lambert W-function; expansion of that W-function
introduces the iterated logarithms
\[
\log\log Z,\qquad \log\log\log Z,
\]
and the lower-order terms of \(\log K\) generate additional algebraic sectors
in inverse powers of the Lambert-normalized variable.  These scales are
strictly ordered but not by a single small parameter.  That is precisely the
setting in which a transseries is the natural language.

\subsection{The central strategy}

The development has four layers.

\begin{enumerate}
\item Center the independent variable at \(x=\tfrac12\).  This eliminates an
otherwise troublesome \(r\log r\) term and leaves an even inverse-power
expansion.
\item Invert the dominant phase exactly with \(W_0\).  This produces a large
reference scale \(\rho\) and a slow variable \(L=\log\rho\).
\item Solve the remaining equation in the complete ring of formal powers of
\(q=\rho^{-2}\), with coefficients depending slowly on \(L\).  The linearized
coefficient is \(L\), so the recurrence is triangular.
\item Improve the coordinate further by introducing
\(T=x^2-x+\tfrac16\).  This absorbs the logarithmic correction into the
leading phase and delays the first unresolved term by a full asymptotic
sector.
\end{enumerate}

The same mechanism works for a broad class of maps with dominant part
\(a x^p(\log x-b)\).  The K-function is especially instructive because both
normalizations can be carried out in closed form and because the Bernoulli
structure of its direct expansion is explicit.

\subsection{Relation to the accompanying transseries program}

The series-and-transseries directory of the ProveIt project is devoted to the
algebra, composition, and reversion of asymptotic expansions rather than to a
single special function \cite{ProveItSeries}.  The present article follows
that operational viewpoint: the K-function supplies a concrete test case,
while the later sections isolate the reusable inversion calculus.  In
particular, exact inversion of the dominant monomial is performed before any
formal reversion, and analytic error estimates are separated from formal
coefficient algebra.

\subsection{Status and scope}

The positive real inverse is treated completely to all algebraic orders.  A
sectorial complex version is also given, but a global classification of the
Riemann surface of the inverse K-function is not attempted.  Exponentially
small terms are discussed through a general transport theorem and through the
known exponentially improved theory of the Barnes G-function
\cite{Nemes2014}; explicit Stokes multipliers for every inverse branch remain
a separate problem.

The coefficient recurrences, the centered inverse expansion, and the
accelerated \(T\)-coordinate expansion were derived for this report and were
checked symbolically and numerically.  No priority claim is made: the
literature on hyperfactorial asymptotics is extensive, including refinements
such as \cite{Xu2020}, and an exhaustive historical search is beyond the
scope of this article.

\section{The K-function and its principal inverse}
\label{sec:k-definition}

\subsection{Definitions and elementary identities}

For positive real \(x\), one convenient definition is
\begin{equation}
K(x)=(2\pi)^{-(x-1)/2}
\exp\left\{\binom{x}{2}+
\int_0^{x-1}\log\Gamma(t+1)\,\dd t\right\}.
\label{eq:K-integral-definition}
\end{equation}
Equivalent complex definitions may be written in terms of the Hurwitz zeta
function or the Barnes G-function:
\begin{equation}
K(z)=\exp\!\left(\zeta'(-1,z)-\zeta'(-1)\right),
\label{eq:K-hurwitz}
\end{equation}
and, on compatible logarithmic branches,
\begin{equation}
K(z)G(z)=\Gamma(z)^{z-1}.
\label{eq:K-Barnes-relation}
\end{equation}
These standard identities are reviewed in \cite{KFunctionWiki,DLMFBarnesG}.
The functional equation is
\begin{equation}
K(z+1)=z^zK(z),
\label{eq:K-functional}
\end{equation}
so that
\begin{equation}
K(n+1)=\prod_{j=1}^{n}j^j.
\label{eq:K-integers}
\end{equation}

Throughout the positive-real analysis we write
\[
\kappa(x)=\log K(x),
\]
with the ordinary real logarithm.

\begin{lemma}[Logarithmic derivatives]
For \(x>0\),
\begin{align}
\kappa'(x)
 &=\log\Gamma(x)+x-\frac12-\frac12\log(2\pi),
\label{eq:kappa-prime}\\
\kappa''(x)&=\psi(x)+1.
\label{eq:kappa-second}
\end{align}
\end{lemma}

\begin{proof}
Differentiate \cref{eq:K-integral-definition}.  The derivative of the upper
limit contributes \(\log\Gamma(x)\), the quadratic term contributes
\(x-\tfrac12\), and the prefactor contributes
\(-\tfrac12\log(2\pi)\).  Differentiating once more gives
\cref{eq:kappa-second}.
\end{proof}

\begin{theorem}[The principal positive inverse]
\label{thm:positive-inverse}
The restriction
\[
K_+=K|_{[2,\infty)}:[2,\infty)\longrightarrow[1,\infty)
\]
is strictly increasing and onto.  Hence it possesses a unique inverse
\[
\Kinv:[1,\infty)\longrightarrow[2,\infty).
\]
\end{theorem}

\begin{proof}
Since \(\psi\) is increasing on \((0,\infty)\) and
\(\psi(2)=1-\gamma>0\), \cref{eq:kappa-second} gives
\(\kappa''(x)>0\) for \(x\ge2\).  Also
\[
\kappa'(2)=\frac32-\frac12\log(2\pi)>0.
\]
Thus \(\kappa'\), and therefore \(K'\), is positive on \([2,\infty)\).
\Cref{eq:K-integers} gives \(K(2)=1\), while the asymptotic formula
proved below implies \(K(x)\to\infty\).  The conclusion follows.
\end{proof}

\begin{remark}
The K-function is not monotone on all of \((0,\infty)\); the restriction in
\cref{thm:positive-inverse} selects the branch that tends to infinity and is
the branch relevant to large arguments.
\end{remark}

\subsection{The Glaisher--Kinkelin constant}

The constant appearing in the asymptotic expansion is
\begin{equation}
A=\exp\!\left(\frac1{12}-\zeta'(-1)\right)
 =1.2824271291006226368\ldots,
\label{eq:glaisher}
\end{equation}
and we shall repeatedly use
\begin{equation}
C=\log A-\frac18
 =0.12375447703378426255\ldots.
\label{eq:C}
\end{equation}

\section{The direct large-argument expansion of \texorpdfstring{$\log K$}{log K}}
\label{sec:direct-asymptotic}

\subsection{An all-orders formula}

We use Bernoulli numbers in the convention
\[
\frac{t}{\e^t-1}=\sum_{n=0}^{\infty}B_n\frac{t^n}{n!}.
\]

\begin{theorem}[Direct K-function expansion]
\label{thm:direct-K-expansion}
For every fixed integer \(N\ge1\), as \(x\to+\infty\),
\begin{align}
\kappa(x)
={}&\left(\frac{x^2}{2}-\frac{x}{2}+\frac1{12}\right)\log x
-\frac{x^2}{4}+\log A
\notag\\
&-\sum_{n=1}^{N-1}
\frac{B_{2n+2}}{(2n)(2n+1)(2n+2)x^{2n}}
+\bigO(x^{-2N}).
\label{eq:direct-K-expansion}
\end{align}
The same expansion holds uniformly in closed subsectors
\(|\arg x|\le\pi-\delta\), with a fixed branch of \(\Log x\).
\end{theorem}

\begin{proof}
From \cref{eq:kappa-prime} and Stirling's expansion,
\begin{align*}
\log\Gamma(x)
={}&\left(x-\frac12\right)\log x-x+\frac12\log(2\pi)
+\sum_{m=1}^{N}
\frac{B_{2m}}{(2m)(2m-1)x^{2m-1}}
+\bigO(x^{-2N-1}),
\end{align*}
we obtain
\begin{equation}
\kappa'(x)=\left(x-\frac12\right)\log x-\frac12
+\sum_{m=1}^{N}
\frac{B_{2m}}{(2m)(2m-1)x^{2m-1}}
+\bigO(x^{-2N-1}).
\label{eq:kappa-prime-stirling}
\end{equation}
Termwise integration yields \cref{eq:direct-K-expansion}, apart from an
integration constant.  That constant is \(\log A\), as follows either from
the classical hyperfactorial asymptotic at integers or from the Barnes
G-function relation \cref{eq:K-Barnes-relation} and the standard asymptotic
normalization of \(G\); see \cite{DLMFBarnesG}.  Uniform sectorial remainder
bounds follow from the corresponding sectorial Stirling expansion and
integration along a ray inside the sector.
\end{proof}

The first terms are
\begin{align}
\kappa(x)\sim{}&
\left(\frac{x^2}{2}-\frac{x}{2}+\frac1{12}\right)\log x
-\frac{x^2}{4}+\log A
+\frac{1}{720x^2}-\frac{1}{5040x^4}
\notag\\
&+\frac{1}{10080x^6}-\frac{1}{9504x^8}
+\frac{691}{3603600x^{10}}-\cdots.
\label{eq:direct-first-terms}
\end{align}

\subsection{Divergence and the least-term scale}

The Bernoulli numbers satisfy
\[
|B_{2n}|\sim \frac{2(2n)!}{(2\pi)^{2n}}.
\]
Consequently, \cref{eq:direct-K-expansion} is normally divergent.  Its
smallest term occurs when the index is proportional to \(\pi x\), and the
optimally truncated remainder has an exponential scale comparable to
\(\e^{-2\pi x}\).  This does not invalidate the algebraic expansion: each
fixed finite truncation is a valid Poincare approximation.  It does mean that
an exponentially improved analysis must eventually append additional
sectorial data, a point revisited in \cref{sec:exponential-sectors}.

\section{The unique centered normal form}
\label{sec:centered}

\subsection{Why the shift by one half is forced}

Let \(x=r+h\) in the leading part of \cref{eq:direct-K-expansion}.  The
coefficient of \(r\log r\) is \(h-\tfrac12\).  Hence
\[
h=\frac12
\]
is the unique affine shift that removes this term.  This is more than a
cosmetic simplification: after the shift, all inverse-power corrections are
even, and the inverse series advances in steps of \(r^{-2}\).

Put
\begin{equation}
r=x-\frac12.
\label{eq:r-definition}
\end{equation}
We use Bernoulli polynomials defined by
\[
\frac{t\e^{ut}}{\e^t-1}
=\sum_{n=0}^{\infty}B_n(u)\frac{t^n}{n!}.
\]
Recall
\[
B_n\!\left(\frac12\right)=\left(2^{1-n}-1\right)B_n.
\]

\begin{theorem}[Centered even expansion]
\label{thm:centered-expansion}
As \(r\to+\infty\),
\begin{equation}
\kappa\!\left(r+\frac12\right)
\sim \left(\frac{r^2}{2}-\frac1{24}\right)\log r
-\frac{r^2}{4}+C+
\sum_{n=1}^{\infty}c_n r^{-2n},
\label{eq:centered-expansion}
\end{equation}
where
\begin{equation}
c_n=-\frac{B_{2n+2}(1/2)}{(2n)(2n+1)(2n+2)}.
\label{eq:c-n}
\end{equation}
For every fixed \(N\), truncation after \(c_{N-1}r^{-2N+2}\) has remainder
\(\bigO(r^{-2N})\).
\end{theorem}

\begin{proof}
The shifted Stirling expansion is
\begin{equation}
\log\Gamma\!\left(r+\frac12\right)
\sim r\log r-r+\frac12\log(2\pi)
+\sum_{m=1}^{\infty}
\frac{B_{2m}(1/2)}{(2m)(2m-1)r^{2m-1}}.
\label{eq:shifted-stirling}
\end{equation}
Substitution into \cref{eq:kappa-prime} gives
\[
\frac{\dd}{\dd r}\kappa\!\left(r+\frac12\right)
\sim r\log r-\frac{1}{24r}
+\sum_{n=1}^{\infty}
\frac{B_{2n+2}(1/2)}{(2n+1)(2n+2)r^{2n+1}}.
\]
Integrating term by term yields \cref{eq:centered-expansion} and
\cref{eq:c-n}.  Substitution of \(x=r+\tfrac12\) into
\cref{eq:direct-K-expansion} fixes the constant as
\(C=\log A-\tfrac18\).
\end{proof}

The first coefficients are
\begin{equation}
\begin{aligned}
c_1&=-\frac7{5760},&
 c_2&=\frac{31}{161280},&
 c_3&=-\frac{127}{1290240},\\
c_4&=\frac{511}{4866048},&
 c_5&=-\frac{1414477}{7380172800},&
 c_6&=\frac{8191}{15335424}.
\end{aligned}
\label{eq:c-first}
\end{equation}

\subsection{Separation of the dominant phase}

Define
\begin{equation}
\Phi(r)=\frac12r^2\log r-\frac14r^2.
\label{eq:Phi}
\end{equation}
Then \cref{eq:centered-expansion} becomes
\begin{equation}
\kappa\!\left(r+\frac12\right)
\sim \Phi(r)-\frac1{24}\log r+C+
\sum_{n\ge1}c_n r^{-2n}.
\label{eq:centered-Phi}
\end{equation}
This form makes the hierarchy transparent:
\[
\Phi(r)\asymp r^2\log r,
\qquad -\frac1{24}\log r+C=\bigO(\log r),
\qquad c_n r^{-2n}=\bigO(r^{-2n}).
\]

\section{Exact inversion of the dominant phase}
\label{sec:dominant-inversion}

\subsection{Lambert W appears without approximation}

Let
\begin{equation}
Y=\log Z.
\label{eq:Y}
\end{equation}
The dominant equation is
\begin{equation}
\Phi(\rho)=Y.
\label{eq:dominant-equation}
\end{equation}
Set \(w=2\log\rho-1\).  Since \(\rho^2=\e^{w+1}\),
\[
Y=\frac14\rho^2(2\log\rho-1)
=\frac{\e}{4}w\e^w.
\]
Thus
\begin{equation}
w=W_0\!\left(\frac{4Y}{\e}\right),
\label{eq:w}
\end{equation}
and
\begin{equation}
\rho=\exp\!\left(\frac{w+1}{2}\right)
=2\sqrt{\frac{Y}{w}},
\qquad
L=\log\rho=\frac{w+1}{2}.
\label{eq:rho-L}
\end{equation}
No asymptotic approximation has yet been made: \(\rho\) is the exact positive
solution of \cref{eq:dominant-equation}.

The zeroth approximation to the inverse K-function is therefore
\begin{equation}
\Kinv(Z)=\frac12+\rho+\smallo(\rho).
\label{eq:leading-inverse}
\end{equation}

\subsection{Elementary iterated-log expansion}

The W-adapted formula \cref{eq:rho-L} is the cleanest representation.  If an
expansion using only elementary logarithms is desired, set
\begin{equation}
U=\log\!\left(\frac{4Y}{\e}\right),
\qquad V=\log U.
\label{eq:U-V}
\end{equation}
The standard large-argument expansion of \(W_0\) \cite{Corless1996,DLMFW}
gives
\begin{align}
\rho
={}&2\sqrt{\frac{Y}{U}}\Bigg[
1+\frac{V}{2U}
+\frac{V(3V-4)}{8U^2}
+\frac{V(5V^2-16V+8)}{16U^3}
\notag\\
&\hspace{31mm}
+\frac{V(105V^3-568V^2+720V-192)}{384U^4}
+\bigO\!\left(\frac{V^5}{U^5}\right)
\Bigg].
\label{eq:rho-elementary}
\end{align}
Thus the leading sector alone contains arbitrarily deep powers of
\(1/\log Y\), with polynomial coefficients in \(\log\log Y\).

\begin{warningbox}
The additive constant \(\tfrac12\) in \(K_+^{-1}(Z)\) is smaller than every
fixed truncation error of the elementary expansion \cref{eq:rho-elementary}:
for each fixed \(N\),
\[
\frac{\rho}{U^N}\longrightarrow\infty.
\]
Hence the correct transseries order is:
\[
\text{complete Lambert/iterated-log sector of size }\rho,
\quad \frac12,
\quad \rho^{-1},\quad \rho^{-3},\ldots.
\]
Writing a few inverse-log terms and then appending \(\tfrac12\) does not
produce an asymptotically ordered one-dimensional series.
\end{warningbox}

\section{The canonical centered inverse transseries}
\label{sec:canonical}

\subsection{The normalized implicit equation}

Put
\begin{equation}
q=\rho^{-2},
\qquad
r=\rho(1+v),
\qquad
v=\sum_{n\ge1}V_n(L)q^n.
\label{eq:q-v}
\end{equation}
The exact dominant identity \(\Phi(\rho)=Y\) suggests normalizing the
remaining equation by \(\rho^2\).  Define
\begin{align}
\mathcal D_L(v)
={}&\frac12(1+v)^2\bigl(L+\log(1+v)\bigr)
-\frac14(1+v)^2-\frac12L+\frac14.
\label{eq:D-L}
\end{align}
Then
\begin{equation}
\mathcal D_L(v)
=Lv+\frac{L+1}{2}v^2+
\sum_{k=3}^{\infty}
\frac{(-1)^{k+1}}{k(k-1)(k-2)}v^k.
\label{eq:D-expansion}
\end{equation}
Let
\begin{equation}
E=E(L)=C-\frac{L}{24}.
\label{eq:E}
\end{equation}
Substitution of \cref{eq:q-v} into \cref{eq:centered-Phi} gives the formal
implicit equation
\begin{equation}
\boxed{
\mathcal D_L(v)+q\left[
E-\frac1{24}\log(1+v)+
\sum_{m\ge1}c_m q^m(1+v)^{-2m}
\right]=0.
}
\label{eq:canonical-implicit}
\end{equation}

\subsection{Formal existence, uniqueness, and recurrence}

Let
\[
\mathscr R=\mathbb Q[C,L,L^{-1}]
\]
and consider the complete \(q\)-adic ring \(\mathscr R[[q]]\).  The variable
\(L\) is treated as a slow coefficient, while \(q\) records the genuinely
small algebraic scale.

\begin{theorem}[Canonical all-orders recurrence]
\label{thm:canonical-recurrence}
\Cref{eq:canonical-implicit} has a unique solution
\[
v\in q\mathscr R[[q]].
\]
Writing
\[
v=\sum_{n\ge1}V_n(L)q^n,
\]
the coefficients are obtained recursively as follows.  Let
\[
v_{<n}=\sum_{j=1}^{n-1}V_j(L)q^j.
\]
Then
\begin{align}
V_n(L)=-\frac1L\coeff{q}{n}\Bigg\{
&\mathcal D_L(v_{<n})
+q\Bigg[E-\frac1{24}\log(1+v_{<n})
\notag\\
&\hspace{20mm}+
\sum_{m\ge1}c_mq^m(1+v_{<n})^{-2m}\Bigg]
\Bigg\}.
\label{eq:V-recurrence}
\end{align}
Only \(c_1,\ldots,c_{n-1}\) can contribute to \(V_n\).
\end{theorem}

\begin{proof}
Let the left side of \cref{eq:canonical-implicit} be
\(\mathcal F(v,q)\).  Then
\[
\mathcal F(0,0)=0,
\qquad
\partial_v\mathcal F(0,0)=L.
\]
Since \(L\) is a unit in \(\mathscr R\), the formal implicit-function theorem
in the complete \(q\)-adic ring gives a unique solution in
\(q\mathscr R[[q]]\).  More concretely, suppose
\(V_1,\ldots,V_{n-1}\) are known.  At order \(q^n\), the only occurrence of
the new coefficient is the linear term \(LV_nq^n\) from
\(\mathcal D_L(v)\).  All other contributions depend only on earlier
coefficients.  Solving the coefficient equation gives
\cref{eq:V-recurrence}.
\end{proof}

\subsection{The first coefficients}

The first three coefficients are
\begin{equation}
V_1=-\frac{E}{L}=\frac1{24}-\frac{C}{L},
\label{eq:V1}
\end{equation}
\begin{equation}
V_2=-\frac1L\left(
\frac{L+1}{2}V_1^2-\frac1{24}V_1+c_1
\right),
\qquad c_1=-\frac7{5760},
\label{eq:V2}
\end{equation}
and
\begin{equation}
V_3=-\frac1L\left(
(L+1)V_1V_2+\frac16V_1^3
-\frac1{24}\left(V_2-\frac12V_1^2\right)
-2c_1V_1+c_2
\right),
\label{eq:V3}
\end{equation}
where \(c_2=31/161280\).  An equivalent closed expression for \(V_2\) is
\begin{equation}
V_2=
\frac{7L^2-240EL-2880(L+1)E^2}{5760L^3}.
\label{eq:V2-E}
\end{equation}
For reference, the next recursion is
\begin{align}
V_4=-\frac1L\Bigg[{}
&(L+1)\left(V_1V_3+\frac12V_2^2\right)
+\frac12V_1^2V_2-\frac1{24}V_1^4
\notag\\
&-\frac1{24}\left(V_3-V_1V_2+\frac13V_1^3\right)
+c_1(-2V_2+3V_1^2)-4c_2V_1+c_3
\Bigg].
\label{eq:V4}
\end{align}
As \(L\to\infty\), the coefficients remain bounded; for example,
\begin{align*}
V_1&=\frac1{24}-\frac{C}{L},\\
V_2&=-\frac1{1152}+\frac{C/24+1/480}{L}
-\frac{C^2}{2L^2}+\bigO(L^{-3}),\\
V_3&=\frac1{27648}-\frac{C/384+311/725760}{L}
+\bigO(L^{-2}).
\end{align*}

\subsection{The inverse expansion and its analytic meaning}

\begin{theorem}[Principal inverse transseries]
\label{thm:principal-inverse-transseries}
Let \(Z\to+\infty\), and define \(Y,w,\rho,L\) by
\cref{eq:Y,eq:w,eq:rho-L}.  Then
\begin{equation}
\boxed{
\Kinv(Z)\sim \frac12+\rho+
\sum_{n=1}^{\infty}V_n(L)\rho^{1-2n},
}
\label{eq:principal-inverse-transseries}
\end{equation}
where the \(V_n\) are given by \cref{eq:V-recurrence}.  More precisely, for
every fixed \(N\ge1\),
\begin{equation}
\Kinv(Z)-\left[
\frac12+\rho+
\sum_{n=1}^{N-1}V_n(L)\rho^{1-2n}
\right]
=\bigO\!\left(\rho^{1-2N}\right).
\label{eq:canonical-error}
\end{equation}
\end{theorem}

\begin{proof}
Formal substitution into \cref{eq:centered-Phi} gives a residual whose
normalized expansion vanishes through the prescribed power of
\(q=\rho^{-2}\).  The direct asymptotic expansion of \(\kappa\) may be taken
to sufficiently many terms to make its own remainder smaller than the first
omitted inverse term.  On the positive ray,
\begin{equation}
\kappa'\!\left(r+\frac12\right)
=r\log r+\bigO(r^{-1})
\asymp \rho L.
\label{eq:kappa-derivative-scale}
\end{equation}
The mean-value theorem therefore converts the residual into an error in
\(r\) one derivative-scale smaller.  The first omitted formal term is
\(V_N(L)\rho^{1-2N}\), and \(V_N(L)=\bigO(1)\) for fixed \(N\).  This proves
\cref{eq:canonical-error}.
\end{proof}

\begin{practicalbox}
The first three useful approximants are
\begin{align*}
X_0(Z)&=\frac12+\rho,\\
X_1(Z)&=X_0(Z)+\frac{V_1(L)}{\rho},\\
X_2(Z)&=X_1(Z)+\frac{V_2(L)}{\rho^3},\\
X_3(Z)&=X_2(Z)+\frac{V_3(L)}{\rho^5}.
\end{align*}
They require only \(W_0\), elementary arithmetic, and the constant \(C\).
One Newton correction using the exact K-function generally improves the
accuracy far beyond the displayed truncation.
\end{practicalbox}

\section{A faster quadratic normal form}
\label{sec:accelerated}

\subsection{A general logarithmic-absorption lemma}

The centered K expansion contains the slow correction
\(-\tfrac1{24}\log r\).  It can be absorbed into a new algebraic coordinate.
The following elementary observation explains why.

\begin{lemma}[Absorbing a logarithmic correction]
\label{lem:absorb-log}
Suppose
\begin{equation}
F(r)=a r^p(\log r-b)+\beta\log r+\gamma+\bigO(r^{-p}),
\qquad a\ne0,
\label{eq:absorb-assumption}
\end{equation}
and set
\begin{equation}
T=r^p+\alpha,
\qquad
\Psi(T)=\frac{a}{p}T(\log T-pb).
\label{eq:general-T-Psi}
\end{equation}
Then
\begin{equation}
\Psi(r^p+\alpha)
=a r^p(\log r-b)+a\alpha\log r
+a\alpha\left(\frac1p-b\right)+\bigO(r^{-p}).
\label{eq:absorb-expansion}
\end{equation}
Thus the unique constant shift that absorbs \(\beta\log r\) is
\begin{equation}
\alpha=\frac{\beta}{a}.
\label{eq:alpha-general}
\end{equation}
\end{lemma}

\begin{proof}
Insert \(T=r^p+\alpha\) into \cref{eq:general-T-Psi} and use
\[
\log T=p\log r+\log(1+\alpha r^{-p})
=p\log r+\alpha r^{-p}+\bigO(r^{-2p}).
\]
Collecting the terms of orders \(r^p\log r\), \(r^p\), \(\log r\), and
\(1\) gives \cref{eq:absorb-expansion}.
\end{proof}

For the K-function,
\[
a=\frac12,
\qquad p=2,
\qquad b=\frac12,
\qquad \beta=-\frac1{24}.
\]
Hence \(\alpha=-\tfrac1{12}\), and the natural coordinate is
\begin{equation}
T=r^2-\frac1{12}=x^2-x+\frac16.
\label{eq:T-K}
\end{equation}
The positive branch is recovered from
\begin{equation}
x=\frac12+\sqrt{T+\frac1{12}}.
\label{eq:x-from-T}
\end{equation}

\subsection{The accelerated direct expansion}

Define
\begin{equation}
\Psi(T)=\frac14T(\log T-1).
\label{eq:Psi-K}
\end{equation}

\begin{theorem}[Quadratic normal form]
\label{thm:T-normal-form}
As \(T\to+\infty\),
\begin{equation}
\kappa(x)\sim \Psi(T)+C+
\sum_{n=1}^{\infty}d_nT^{-n},
\qquad T=x^2-x+\frac16,
\label{eq:T-normal-form}
\end{equation}
where, with \(a_0=1/12\),
\begin{equation}
d_n=
\frac{(-1)^n a_0^{n+1}}{4(n+1)}
+\sum_{m=1}^{n}
(-1)^{n-m}\binom{n-1}{m-1}a_0^{n-m}c_m.
\label{eq:d-n}
\end{equation}
The first coefficients are
\begin{equation}
\begin{aligned}
d_1&=-\frac1{480},&
 d_2&=\frac{31}{90720},&
 d_3&=-\frac{103}{725760},\\
d_4&=\frac{179}{1330560},&
 d_5&=-\frac{6480899}{28021593600},&
 d_6&=\frac{3485299}{5604318720}.
\end{aligned}
\label{eq:d-first}
\end{equation}
\end{theorem}

\begin{proof}
Because \(r^2=T+a_0\) and
\(\tfrac12r^2-\tfrac1{24}=T/2\), the first two terms of
\cref{eq:centered-expansion} satisfy the exact identity
\begin{align*}
\left(\frac{r^2}{2}-\frac1{24}\right)\log r-\frac{r^2}{4}
&=\frac{T}{4}\log(T+a_0)-\frac{T}{4}-\frac{a_0}{4}\\
&=\Psi(T)+\frac{T}{4}\log\left(1+\frac{a_0}{T}\right)-\frac{a_0}{4}.
\end{align*}
The term with first power of \(a_0/T\) cancels \(a_0/4\).  Expanding the
remaining logarithm and
\[
(T+a_0)^{-m}=T^{-m}(1+a_0/T)^{-m}
\]
then collecting the coefficient of \(T^{-n}\) gives \cref{eq:d-n}.
\end{proof}

The importance of \cref{eq:T-normal-form} is structural.  The correction of
size \(\log r\), which produced a displacement of size \(r^{-1}\) in the
canonical inverse, has disappeared.  The first unresolved direct term is now
\(d_1/T\), and it produces an inverse displacement only of order
\(T^{-3/2}/\log T\).

\subsection{Exact leading inversion in the T-coordinate}

Set
\begin{equation}
\eta=Y-C.
\label{eq:eta}
\end{equation}
The dominant equation \(\Psi(\tau)=\eta\) is
\[
\frac14\tau(\log\tau-1)=\eta.
\]
Let
\begin{equation}
\omega=W_0\!\left(\frac{4\eta}{\e}\right),
\qquad
\tau=\frac{4\eta}{\omega}=\e^{\omega+1},
\qquad
S=\log\tau=\omega+1.
\label{eq:tau-S}
\end{equation}
Then \(\Psi(\tau)=\eta\) exactly.  The leading accelerated approximation is
therefore
\begin{equation}
\boxed{
X_{\mathrm{acc},0}(Z)
=\frac12+
\sqrt{\frac1{12}+
\frac{4(\log Z-C)}{W_0(4(\log Z-C)/\e)}}.
}
\label{eq:accelerated-leading}
\end{equation}

\subsection{All-orders accelerated recurrence}

Write
\begin{equation}
T=\tau(1+u),
\qquad
\nu=\tau^{-1},
\qquad
\nu\to0^+.
\label{eq:T-u}
\end{equation}
Define
\begin{equation}
\mathcal E_S(u)
=(1+u)(S+\log(1+u)-1)-(S-1).
\label{eq:E-S}
\end{equation}
Its expansion is
\begin{equation}
\mathcal E_S(u)=Su+
\sum_{k=2}^{\infty}\frac{(-1)^k}{k(k-1)}u^k.
\label{eq:E-S-expansion}
\end{equation}
After multiplying the equation by \(4/\tau\), the normalized implicit
relation becomes
\begin{equation}
\boxed{
\mathcal E_S(u)+4\sum_{m\ge1}d_m\nu^{m+1}(1+u)^{-m}=0.
}
\label{eq:accelerated-implicit}
\end{equation}
Since the perturbation starts at \(\nu^2\), so does \(u\).

\begin{theorem}[Accelerated all-orders recurrence]
\label{thm:accelerated-recurrence}
\Cref{eq:accelerated-implicit} has a unique solution
\begin{equation}
u=\sum_{n=2}^{\infty}U_n(S)\nu^n
\quad\text{in}\quad
\nu^2\mathbb Q[S,S^{-1}][[\nu]].
\label{eq:u-series}
\end{equation}
For
\[
u_{<n}=\sum_{j=2}^{n-1}U_j(S)\nu^j,
\]
the triangular recurrence is
\begin{equation}
U_n(S)=-\frac1S\coeff{\nu}{n}\left\{
\mathcal E_S(u_{<n})
+4\sum_{m\ge1}d_m\nu^{m+1}(1+u_{<n})^{-m}
\right\}.
\label{eq:U-recurrence}
\end{equation}
The first coefficients are
\begin{align}
U_2&=\frac1{120S},
&U_3&=-\frac{31}{22680S},
\label{eq:U2U3}\\
U_4&=\frac{1030S^2-126S-63}{1814400S^3},
&U_5&=-\frac{16110S^2-1023S-341}{29937600S^3}.
\label{eq:U4U5}
\end{align}
\end{theorem}

\begin{proof}
The left side of \cref{eq:accelerated-implicit} vanishes at
\((u,\nu)=(0,0)\), and its derivative with respect to \(u\) there is \(S\),
a unit in the chosen coefficient ring.  Formal implicit inversion therefore
gives a unique series.  Since the forcing term begins with
\(4d_1\nu^2\), the coefficients of \(\nu^0\) and \(\nu^1\) vanish and
\(u=\bigO(\nu^2)\).  Coefficient extraction gives
\cref{eq:U-recurrence}.  Direct calculation gives
\cref{eq:U2U3,eq:U4U5}.
\end{proof}

\begin{corollary}[Accelerated inverse transseries]
\label{cor:accelerated-inverse}
As \(Z\to+\infty\),
\begin{equation}
\boxed{
\Kinv(Z)\sim
\frac12+
\sqrt{\frac1{12}+
\tau\left(1+\sum_{n=2}^{\infty}U_n(S)\tau^{-n}\right)},
}
\label{eq:accelerated-inverse}
\end{equation}
where \(\tau,S\) are given by \cref{eq:tau-S}.  If the sum is truncated after
\(U_N\tau^{-N}\), the error in \(x\) is
\begin{equation}
\bigO\!\left(\tau^{-N-1/2}\right)
\qquad (N\ge2),
\label{eq:accelerated-error}
\end{equation}
with slow powers of \(S^{-1}\) included in the coefficient.
\end{corollary}

The first two corrections may be displayed without expanding the square
root.  If
\[
R=\sqrt{\tau+\frac1{12}},
\]
then
\begin{align}
\Kinv(Z)
={}&\frac12+R
+\frac{1}{240S\tau R}
-\frac{31}{45360S\tau^2R}
+\bigO\!\left(\tau^{-7/2}\right).
\label{eq:accelerated-expanded-x}
\end{align}

\subsection{The continuous hyperfactorial inverse}

Define the continuous hyperfactorial interpolation
\[
\mathcal H(t)=K(t+1),
\qquad t\ge1.
\]
Then
\begin{equation}
\mathcal H^{-1}(Z)=\Kinv(Z)-1.
\label{eq:H-inverse-shift}
\end{equation}
In the accelerated coordinate,
\[
T=t^2+t+\frac16,
\]
so \cref{eq:accelerated-inverse} immediately gives
\begin{equation}
\mathcal H^{-1}(Z)\sim
-\frac12+
\sqrt{\frac1{12}+
\tau\left(1+\sum_{n=2}^{\infty}U_n(S)\tau^{-n}\right)}.
\label{eq:H-inverse}
\end{equation}
This explains the naturally occurring quadratic combination
\(t^2+t+\tfrac16\) in refined hyperfactorial asymptotics.

\section{Numerical verification}
\label{sec:numerics}

\subsection{Reference computation}

For positive real \(x\), the Barnes relation gives the stable identity
\begin{equation}
\kappa(x)=(x-1)\log\Gamma(x)-\log G(x).
\label{eq:numeric-logK}
\end{equation}
Reference inverse values were obtained by Newton iteration
\begin{equation}
x_{j+1}=x_j-
\frac{\kappa(x_j)-Y}
{\log\Gamma(x_j)+x_j-\tfrac12-\tfrac12\log(2\pi)},
\label{eq:newton-K}
\end{equation}
using high-precision implementations of \(\Gamma\), \(G\), and \(W_0\).
The signed errors below are approximation minus reference value.

\subsection{Canonical centered expansion}

Let \(X_j\) denote the canonical approximant through \(V_j\), with
\(X_0=\tfrac12+\rho\).  The results are:

\begin{table}[htbp]
\centering
\scriptsize
\caption{Canonical centered inverse: signed errors.}
\label{tab:canonical-errors}
\begin{tabular}{@{}r r r r r r@{}}
\toprule
\(Y\) & reference \(x\) & \(X_0-x\) & \(X_1-x\) & \(X_2-x\) & \(X_3-x\)\\
\midrule
10
& 4.966157945496709
& \(9.1616267\,10^{-3}\)
& \(1.9348166\,10^{-5}\)
& \(1.3471239\,10^{-7}\)
& \(-6.0356982\,10^{-10}\)\\
100
& 10.915136726993085
& \(1.0696633\,10^{-3}\)
& \(-2.0464973\,10^{-7}\)
& \(5.1650691\,10^{-10}\)
& \(-2.6895791\,10^{-12}\)\\
1000
& 27.284589040600087
& \(-1.5034726\,10^{-4}\)
& \(-2.1337345\,10^{-8}\)
& \(4.2407143\,10^{-12}\)
& \(-3.0318935\,10^{-15}\)\\
10000
& 73.181059301767169
& \(-1.7601786\,10^{-4}\)
& \(-7.9953926\,10^{-10}\)
& \(2.7035317\,10^{-14}\)
& \(-2.3761398\,10^{-18}\)\\
\(10^6\)
& 584.235864411692401
& \(-3.8094767\,10^{-5}\)
& \(-2.5238008\,10^{-13}\)
& \(6.1276265\,10^{-19}\)
& \(-8.3193819\,10^{-25}\)\\
\bottomrule
\end{tabular}
\end{table}

\subsection{Accelerated quadratic expansion}

Let \(A_0\) be \cref{eq:accelerated-leading}, and let \(A_j\) for
\(j=2,3,4,5\) include the terms through \(U_j\) inside the square root.

\begin{table}[htbp]
\centering
\scriptsize
\caption{Accelerated inverse: signed errors.}
\label{tab:accelerated-errors}
\begin{tabular}{@{}r r r r r r@{}}
\toprule
\(Y\) & \(A_0-x\) & \(A_2-x\) & \(A_3-x\) & \(A_4-x\) & \(A_5-x\)\\
\midrule
10
& \(-1.5587271\,10^{-5}\)
& \(1.2729709\,10^{-7}\)
& \(-2.4670503\,10^{-9}\)
& \(1.1659381\,10^{-10}\)
& \(-9.8337554\,10^{-12}\)\\
100
& \(-7.8648990\,10^{-7}\)
& \(1.1875374\,10^{-9}\)
& \(-4.3964091\,10^{-12}\)
& \(3.8746388\,10^{-14}\)
& \(-6.1491072\,10^{-16}\)\\
1000
& \(-3.2972642\,10^{-8}\)
& \(7.5368039\,10^{-12}\)
& \(-4.2730563\,10^{-15}\)
& \(5.6974464\,10^{-18}\)
& \(-1.3723490\,10^{-20}\)\\
10000
& \(-1.2659858\,10^{-9}\)
& \(3.9307274\,10^{-14}\)
& \(-3.0434681\,10^{-18}\)
& \(5.5017139\,10^{-22}\)
& \(-1.7988347\,10^{-25}\)\\
\(10^6\)
& \(-1.6444009\,10^{-12}\)
& \(7.9154321\,10^{-19}\)
& \(-9.5514849\,10^{-25}\)
& \(2.6701788\,10^{-30}\)
& \(-1.3514777\,10^{-35}\)\\
\bottomrule
\end{tabular}
\end{table}

The accelerated leading approximation already outperforms the first
corrected canonical formula over most of the table.  This is not accidental:
it has exactly resummed the constant \(C\), the logarithmic correction, and
the algebraic effects induced by the quadratic change of variable.

\section{A general theory of power--logarithmic transseries inversion}
\label{sec:general-theory}

\subsection{The model class}

Let \(p>0\), \(\sigma>0\), and \(a\ne0\).  Consider a sectorial asymptotic
expansion of the form
\begin{equation}
F(x)\sim a x^p(\log x-b)
+\sum_{m=0}^{\infty}
 x^{p-(m+1)\sigma}A_m(\log x).
\label{eq:general-class}
\end{equation}
The functions \(A_m\) belong to a slow coefficient field \(\mathscr S\).
Typical choices are rational functions of \(\log x\), polynomial-logarithmic
transseries in deeper logarithms, or sectorial analytic functions admitting
such expansions.  The essential requirement is that the coefficient field
be closed under differentiation, composition with \(\log(1+v)\), and
inversion of the linearized factor introduced below.

The K-function in the centered variable belongs to this class with
\begin{equation}
 a=\frac12,
 \qquad p=2,
 \qquad b=\frac12,
 \qquad \sigma=2,
\label{eq:K-general-parameters}
\end{equation}
and
\[
A_0(L)=C-\frac{L}{24},
\qquad A_m(L)=c_m\quad(m\ge1).
\]

\subsection{Exact dominant inversion}

\begin{theorem}[Lambert--W normalizer]
\label{thm:general-W-normalizer}
Fix a branch \(W_k\) and compatible branches of the logarithm and
\(p\)-th root.  The equation
\begin{equation}
y=a\rho^p(\log\rho-b)
\label{eq:general-dominant}
\end{equation}
has the solution
\begin{align}
\Omega&=W_k\!\left(\frac{p}{a}\e^{-pb}y\right),
\label{eq:Omega-general}\\
\Lambda&=\log\rho=b+\frac{\Omega}{p},
\label{eq:Lambda-general}\\
\rho&=\exp\!\left(b+\frac{\Omega}{p}\right)
=\left[
\frac{py/a}{W_k((p/a)\e^{-pb}y)}
\right]^{1/p}.
\label{eq:rho-general}
\end{align}
\end{theorem}

\begin{proof}
Put \(\Omega=p(\log\rho-b)\).  Then
\[
\rho^p=\e^{pb+\Omega},
\]
and \cref{eq:general-dominant} becomes
\[
\Omega\e^\Omega=\frac{p}{a}\e^{-pb}y.
\]
Applying the selected W branch gives \cref{eq:Omega-general}; the remaining
formulas follow immediately.
\end{proof}

\subsection{The normalized equation and formal theorem}

Set
\begin{equation}
q=\rho^{-\sigma},
\qquad x=\rho(1+v).
\label{eq:general-q-v}
\end{equation}
Define
\begin{equation}
\mathcal D_{p,\Lambda,b}(v)
=(1+v)^p\bigl(\Lambda-b+\log(1+v)\bigr)
-(\Lambda-b).
\label{eq:general-D}
\end{equation}
After division by \(a\rho^p\), the inversion equation is
\begin{align}
0={}&\mathcal D_{p,\Lambda,b}(v)
+\frac1a\sum_{m\ge0}q^{m+1}
(1+v)^{p-(m+1)\sigma}
\notag\\
&\hspace{35mm}\times
A_m\bigl(\Lambda+\log(1+v)\bigr).
\label{eq:general-implicit}
\end{align}
The linear coefficient is
\begin{equation}
\partial_v\mathcal D_{p,\Lambda,b}(0)
=p(\Lambda-b)+1=\Omega+1.
\label{eq:linear-factor}
\end{equation}

\begin{theorem}[Formal inversion over a slow coefficient field]
\label{thm:general-formal-inversion}
Assume \(\Omega+1\) is invertible in the slow coefficient field
\(\mathscr S\).  Then \cref{eq:general-implicit} has a unique solution
\begin{equation}
v=\sum_{n\ge1}V_n q^n\in q\mathscr S[[q]].
\label{eq:general-v-series}
\end{equation}
If \(v_{<n}=\sum_{j=1}^{n-1}V_jq^j\), then
\begin{align}
V_n=-\frac1{\Omega+1}\coeff{q}{n}\Bigg\{
&\mathcal D_{p,\Lambda,b}(v_{<n})
\notag\\
&+\frac1a\sum_{m\ge0}q^{m+1}
(1+v_{<n})^{p-(m+1)\sigma}
\notag\\
&\hspace{28mm}\times
A_m\bigl(\Lambda+\log(1+v_{<n})\bigr)
\Bigg\}.
\label{eq:general-V-recurrence}
\end{align}
Consequently,
\begin{equation}
F^{-1}(y)\sim
\rho(y)\left(1+\sum_{n\ge1}V_n(y)\rho(y)^{-n\sigma}\right).
\label{eq:general-inverse}
\end{equation}
\end{theorem}

\begin{proof}
The left side of \cref{eq:general-implicit} vanishes at \((v,q)=(0,0)\),
and by \cref{eq:linear-factor} its derivative with respect to \(v\) is the
unit \(\Omega+1\).  The formal implicit-function theorem gives existence and
uniqueness in the complete \(q\)-adic ring.  At order \(q^n\), the new
coefficient occurs only as \((\Omega+1)V_nq^n\); coefficient extraction gives
\cref{eq:general-V-recurrence}.
\end{proof}

\subsection{Universal first coefficients}

Let
\begin{equation}
B=\Lambda-b=\frac{\Omega}{p},
\qquad
D_2=\frac12\left[p(p-1)B+2p-1\right].
\label{eq:general-D2}
\end{equation}
Then
\begin{equation}
V_1=-\frac{A_0(\Lambda)}{a(\Omega+1)},
\label{eq:general-V1}
\end{equation}
and
\begin{align}
V_2=-\frac1{\Omega+1}\Bigg[
&D_2V_1^2
+\frac1a\left(
A_0'(\Lambda)+(p-\sigma)A_0(\Lambda)
\right)V_1
\notag\\
&+\frac1aA_1(\Lambda)
\Bigg].
\label{eq:general-V2}
\end{align}
These formulas show explicitly how three effects combine: curvature of the
dominant phase, variation of the slow coefficient, and the next direct
asymptotic block.

\subsection{Analytic promotion by residual control}

Formal algebra determines coefficients, but asymptotic validity is an
analytic statement.  The following elementary theorem is the bridge.

\begin{theorem}[Residual-to-inverse error]
\label{thm:residual-error}
Let \(F\) be continuously differentiable and strictly monotone on an interval
containing the exact inverse point \(x_*=F^{-1}(y)\) and an approximation
\(x_N\).  Put
\[
R_N=F(x_N)-y.
\]
Then, for some \(\xi\) between \(x_N\) and \(x_*\),
\begin{equation}
x_N-x_*=\frac{R_N}{F'(\xi)}.
\label{eq:residual-error}
\end{equation}
If \(F\) has the leading form in \cref{eq:general-class}, then
\begin{equation}
F'(\rho)\sim a\rho^{p-1}(\Omega+1).
\label{eq:general-derivative-scale}
\end{equation}
Hence a normalized residual of order \(q^{N+1}\) produces an inverse error
of the same order as the first omitted transseries term.
\end{theorem}

\begin{proof}
The mean-value theorem gives
\[
F(x_N)-F(x_*)=F'(\xi)(x_N-x_*),
\]
which is \cref{eq:residual-error}.  Differentiating
\(a x^p(\log x-b)\) gives
\[
a x^{p-1}\bigl(p(\log x-b)+1\bigr),
\]
and evaluation at \(x=\rho\) gives \cref{eq:general-derivative-scale}.
\end{proof}

This theorem supplies a useful certification workflow: derive the formal
coefficients, substitute the truncation into the original asymptotic
expansion, estimate the residual, and divide by the derivative scale.  It
also makes clear why a small residual in the function value can correspond
to an even smaller error in the inverse when \(F'\) is large.

\subsection{Newton acceleration in the transseries ring}

\begin{proposition}[Quadratic improvement]
\label{prop:newton-transseries}
Suppose \(x_0=\rho(1+v_0)\) has normalized error
\(v_0-v=\bigO(q^m)\), with \(m\ge1\), and the linear factor
\(\Omega+1\) is a unit.  One formal Newton step
\begin{equation}
x_1=x_0-\frac{F(x_0)-y}{F'(x_0)}
\label{eq:formal-newton}
\end{equation}
then has normalized error \(v_1-v=\bigO(q^{2m})\).
\end{proposition}

\begin{proof}
Taylor expansion around the exact formal solution gives
\[
F(x_0)-y=F'(x)(x_0-x)+\bigO((x_0-x)^2).
\]
Because the normalized derivative is a unit, division is permitted in the
coefficient ring.  The linear error cancels in \cref{eq:formal-newton},
leaving a quadratic remainder.
\end{proof}

Newton's method is therefore not merely a numerical afterthought.  It is an
algebraic accelerator that can double the number of correct \(q\)-adic
blocks at each step.

\subsection{Nonlattice supports and Hahn-field extension}

The lattice \(q^n=\rho^{-n\sigma}\) is convenient but not essential.  If the
direct expansion contains powers with exponents in a well-ordered additive
semigroup \(\mathcal M\subset\mathbb R_{\ge0}\), one replaces
\(\mathscr S[[q]]\) by the corresponding Hahn series ring
\[
\mathscr S((\mathcal M)).
\]
The same proof works coefficient by coefficient: at every monomial, the new
unknown appears multiplied by \(\Omega+1\), while all remaining terms depend
on earlier monomials.  Nested logarithms or additional exponential scales
are handled by enlarging the slow coefficient field and the monomial group,
provided the support remains well ordered.

\section{Coordinate normal forms beyond the K-function}
\label{sec:normal-forms-general}

\subsection{Why normal forms matter}

The canonical theorem applies directly to \cref{eq:general-class}, but it may
not give the most efficient coordinates.  A lower-order term such as
\(\beta\log x\) creates an inverse displacement much larger than subsequent
algebraic corrections.  Lemma \ref{lem:absorb-log} shows that this entire
sector can be absorbed by an algebraic change of variable
\(T=x^p+\beta/a\).  The resulting leading equation is again inverted by W,
now in \(T\) rather than \(x\).

Suppose
\begin{equation}
F(x)=a x^p(\log x-b)+\beta\log x+\gamma+\bigO(x^{-p}).
\label{eq:normal-form-input}
\end{equation}
Set
\begin{equation}
T=x^p+\frac{\beta}{a},
\qquad
\gamma_*=\gamma-\beta\left(\frac1p-b\right).
\label{eq:normal-form-T-gamma}
\end{equation}
Then
\begin{equation}
F(x)=\frac{a}{p}T(\log T-pb)+\gamma_*+\bigO(T^{-1}).
\label{eq:normal-form-output}
\end{equation}
The exact dominant inverse is
\begin{equation}
T_0(y)=
\frac{p(y-\gamma_*)/a}
{W_k\!\left((p/a)\e^{-pb}(y-\gamma_*)\right)}.
\label{eq:normal-form-inverse}
\end{equation}
For the K-function, \(p=2\), \(a=1/2\), \(b=1/2\), and
\(\gamma_*=C\), reproducing \cref{eq:accelerated-leading}.

\subsection{A normal-form algorithm}

For more complicated direct expansions, one can seek
\begin{equation}
T=x^p+\alpha_0+\alpha_1x^{-\sigma}+\alpha_2x^{-2\sigma}+\cdots
\label{eq:formal-coordinate}
\end{equation}
so that \(F(x)\), expressed as a function of \(T\), contains as few slow
logarithmic corrections as possible.  Coefficients are chosen successively
by matching the unwanted terms.  This is analogous to normal-form reduction
in dynamical systems: a coordinate change removes nonresonant terms, while
terms that cannot be removed become invariants of the asymptotic class.

For K, the single constant shift in \(r^2\) already removes the only
logarithmic correction, and the remaining expansion is a pure inverse-power
series in \(T\).  This unusually clean outcome is the source of the rapid
convergence observed in \cref{tab:accelerated-errors}.

\section{Complex inverse branches}
\label{sec:complex-branches}

\subsection{Two independent branch indices}

Let \(Z\) tend to infinity in a complex sector.  Choose a branch of
\(\Log K(z)\) in a corresponding large-\(z\) sector and a branch of
\(\Log Z\).  The equation \(K(z)=Z\) lifts locally to
\begin{equation}
\kappa(z)=Y_m,
\qquad
Y_m=\Log Z+2\pi\ii m,
\qquad m\in\mathbb Z.
\label{eq:Y-m}
\end{equation}
The integer \(m\) records the logarithmic sheet of the outer exponential.
The dominant phase introduces a second index, the Lambert branch \(W_k\).
Define
\begin{equation}
w_{m,k}=W_k\!\left(\frac{4Y_m}{\e}\right),
\qquad
\rho_{m,k}=\exp\!\left(\frac{w_{m,k}+1}{2}\right),
\qquad
L_{m,k}=\frac{w_{m,k}+1}{2}.
\label{eq:complex-rho}
\end{equation}
Whenever \(\rho_{m,k}\) lies in the selected asymptotic sector and the
logarithms are compatible, the same coefficient recurrence gives
\begin{equation}
z_{m,k}(Z)\sim
\frac12+\rho_{m,k}+
\sum_{n\ge1}V_n(L_{m,k})\rho_{m,k}^{1-2n}.
\label{eq:complex-inverse}
\end{equation}

\begin{remark}
Formula \cref{eq:complex-inverse} is a sectorial branch construction, not a
global enumeration theorem.  Different pairs \((m,k)\) may describe the
same analytic branch after continuation, and some pairs may violate the
chosen logarithmic sector.  A complete quotient by these identifications
requires a separate global study of the inverse Riemann surface.
\end{remark}

\subsection{The resonance or critical case}

In the general theory, ordinary implicit inversion fails when
\begin{equation}
\Omega+1=0.
\label{eq:resonance}
\end{equation}
This is exactly the point where the derivative of the dominant map vanishes.
Since \(\Omega=W_k(\cdot)\), it corresponds to the Lambert branch point
\(-\e^{-1}\).  Near such a point, the correct local inverse scale is
Puiseux-like, beginning with a square root, rather than an ordinary
\(q\)-series.  On the principal positive branch, \(\Omega>0\), so this
resonance never occurs.

\subsection{The accelerated complex form}

The quadratic coordinate also extends sectorially.  With
\begin{equation}
\eta_m=Y_m-C,
\qquad
\omega_{m,k}=W_k\!\left(\frac{4\eta_m}{\e}\right),
\qquad
\tau_{m,k}=\frac{4\eta_m}{\omega_{m,k}},
\label{eq:complex-tau}
\end{equation}
the expansion is
\begin{equation}
z_{m,k}(Z)\sim
\frac12+
\left[
\frac1{12}+\tau_{m,k}
\left(1+\sum_{n\ge2}U_n(S_{m,k})\tau_{m,k}^{-n}\right)
\right]^{1/2},
\label{eq:complex-accelerated}
\end{equation}
where \(S_{m,k}=\Log\tau_{m,k}\), and the square-root branch must be chosen
consistently with the target sector.

\section{Exponentially small sectors and their inversion}
\label{sec:exponential-sectors}

\subsection{What lies beyond the algebraic series}

The coefficients in \cref{eq:centered-expansion} grow factorially.  The
optimally truncated remainder is exponentially small, and the Barnes
G-function has a well-developed exponentially improved asymptotic theory
with Stokes smoothing \cite{Nemes2014}.  Through
\cref{eq:K-Barnes-relation}, the same phenomenon enters the K-function.
The explicit form depends on the complex sector: Stokes lines and terminant
functions cannot be represented by a single globally valid power series.

For inversion, however, the leading transport mechanism is universal.

\begin{theorem}[Transport of a first exponential sector]
\label{thm:exponential-transport}
Suppose, in a fixed sector,
\begin{equation}
F(x)=F_{\mathrm{alg}}(x)+
S(x)\e^{-\lambda x}+\text{smaller exponential sectors},
\qquad \Re(\lambda x)>0,
\label{eq:direct-exponential}
\end{equation}
and let \(x_{\mathrm{alg}}(y)\) solve
\(F_{\mathrm{alg}}(x)=y\).  If
\(F_{\mathrm{alg}}'(x_{\mathrm{alg}})\ne0\), then the inverse has the first
exponential correction
\begin{equation}
F^{-1}(y)=x_{\mathrm{alg}}(y)
-\frac{S(x_{\mathrm{alg}}(y))}
{F_{\mathrm{alg}}'(x_{\mathrm{alg}}(y))}
\e^{-\lambda x_{\mathrm{alg}}(y)}
+\cdots.
\label{eq:inverse-exponential}
\end{equation}
\end{theorem}

\begin{proof}
Write \(x=x_{\mathrm{alg}}+\delta\) and substitute into
\cref{eq:direct-exponential}.  Taylor expansion gives
\[
0=F_{\mathrm{alg}}'(x_{\mathrm{alg}})\delta
+S(x_{\mathrm{alg}})\e^{-\lambda x_{\mathrm{alg}}}
+\text{quadratically or exponentially smaller terms}.
\]
Solving the leading linear equation yields \cref{eq:inverse-exponential}.
\end{proof}

For the K-function, \(F_{\mathrm{alg}}'=\kappa'\sim r\log r\).  Thus a
direct exponential contribution of scale \(\e^{-2\pi r}\) is transported to
an inverse contribution of scale
\begin{equation}
\frac{\e^{-2\pi\rho}}{\rho L}
\label{eq:K-exponential-scale}
\end{equation}
up to its sector-dependent prefactor.  Higher exponential sectors interact
through the usual transseries convolution rules.

\section{Applications to related functions}
\label{sec:applications}

\subsection{Generalized hyperfactorials}

For \(s>-1\), consider the discrete generalized hyperfactorial
\begin{equation}
H_s(n)=\prod_{j=1}^{n}j^{j^s}.
\label{eq:Hs}
\end{equation}
Euler--Maclaurin summation applied to
\(\sum_{j\le n}j^s\log j\) begins with
\begin{equation}
\log H_s(n)
\sim \frac{n^{s+1}}{s+1}\log n
-\frac{n^{s+1}}{(s+1)^2}
+\frac12n^s\log n+\cdots.
\label{eq:Hs-leading}
\end{equation}
Ignoring for the moment the boundary correction and choosing a smooth
interpolation, the dominant parameters are
\begin{equation}
p=s+1,
\qquad a=\frac1{s+1},
\qquad b=\frac1{s+1}.
\label{eq:Hs-parameters}
\end{equation}
Therefore the first Lambert-normalized scale is
\begin{equation}
\rho_s(Y)=
\left[
\frac{(s+1)^2Y}
{W_0((s+1)^2Y/\e)}
\right]^{1/(s+1)}.
\label{eq:Hs-inverse-leading}
\end{equation}
The lower-order Euler--Maclaurin terms determine the centering and the
coefficient functions \(A_m\).  Formula \cref{eq:general-V-recurrence} then
constructs the inverse to all orders.  The ordinary hyperfactorial is
\(s=1\), for which \cref{eq:Hs-inverse-leading} reduces to
\(2\sqrt{Y/W_0(4Y/\e)}\).

\subsection{The inverse gamma scale}

Stirling's formula for \(\log\Gamma(x)\) has dominant part
\[
x(\log x-1),
\]
corresponding to \(a=1\), \(p=1\), \(b=1\).  The leading inverse scale is
therefore
\begin{equation}
\rho(Y)=\frac{Y}{W_0(Y/\e)}.
\label{eq:inverse-gamma-leading}
\end{equation}
The familiar inverse-factorial and inverse-gamma expansions are thus the
\(p=1\) member of the same theory, while the K-function is the \(p=2\)
member.

\subsection{Barnes-type and multiple-gamma functions}

Logarithms of Barnes multiple gamma functions have leading terms of the form
\(x^p\log x\) with integer \(p\).  After removal of polynomial terms and
appropriate centering, their inverse functions fit
\cref{eq:general-class}.  The dominant W formula
\cref{eq:rho-general} is therefore universal across the hierarchy; only the
slow coefficient field and the exponent lattice change.  In higher order,
several polynomial-logarithmic terms may need to be absorbed by a more
elaborate coordinate normal form before the residual recurrence becomes
sparse.

\section{Algorithms and reproducibility}
\label{sec:algorithms}

\subsection{Coefficient-generation algorithm}

The canonical K coefficients can be generated without expanding the entire
expression at once.

\begin{algorithm}[Triangular coefficient generation]
\label{alg:triangular}
Given a target order \(N\):
\begin{enumerate}
\item Compute \(c_1,\ldots,c_{N-1}\) from \cref{eq:c-n}.
\item Initialize \(v=0\).
\item For \(n=1,\ldots,N\), truncate every power series in
\cref{eq:canonical-implicit} modulo \(q^{n+1}\), extract the coefficient of
\(q^n\), divide its known part by \(-L\), and append \(V_nq^n\) to \(v\).
\item Optionally simplify \(V_n\) as a Laurent polynomial in \(L\) and
\(C\).
\item Verify by substituting the resulting \(v\) back into
\cref{eq:canonical-implicit} and checking that the residual is
\(\bigO(q^{N+1})\).
\end{enumerate}
\end{algorithm}

The accelerated coefficients are generated identically from
\cref{eq:accelerated-implicit}, except that the unknown series starts at
\(\nu^2\) and the divisor is \(S\).

\subsection{Stable numerical evaluation}

For very large \(Z\), one should never form \(Z\) explicitly.  The natural
input is \(Y=\log Z\).  Compute \(W_0(4Y/\e)\) or
\(W_0(4(Y-C)/\e)\), evaluate the selected transseries truncation, and only
then use the exact special functions if a Newton correction is desired.
This avoids overflow and respects the asymptotic geometry of the problem.

A residual provides an a posteriori error estimate:
\begin{equation}
\widehat\epsilon=
\frac{\kappa(x_{\mathrm{app}})-Y}{\kappa'(x_{\mathrm{app}})}.
\label{eq:posteriori}
\end{equation}
When the approximation is already in the asymptotic basin,
\(x_{\mathrm{app}}-x_*\approx\widehat\epsilon\), and subtracting
\(\widehat\epsilon\) is one Newton step.

\subsection{Reference Python code}

The following script reproduces the numerical tables.  It uses \texttt{mpmath}
for high-precision gamma, Barnes G, and Lambert W evaluations.

\begin{lstlisting}[language=Python,caption={Numerical verification of the inverse K-function transseries.}]
import mpmath as mp

mp.mp.dps = 80

# log(A) = 1/12 - zeta'(-1), and C = log(A) - 1/8.
logA = mp.mpf(1) / 12 - mp.diff(lambda s: mp.zeta(s), -1)
C = logA - mp.mpf(1) / 8


def logK(x):
    """Real log K(x), using K(x) G(x) = Gamma(x)^(x-1)."""
    x = mp.mpf(x)
    return (x - 1) * mp.loggamma(x) - mp.log(mp.barnesg(x))


def dlogK(x):
    """Derivative of log K(x)."""
    x = mp.mpf(x)
    return (mp.loggamma(x) + x - mp.mpf("0.5")
            - mp.log(2 * mp.pi) / 2)


def inverse_logK(y):
    """Solve log K(x) = y on x >= 2 by Newton iteration."""
    y = mp.mpf(y)
    w = mp.lambertw(4 * y / mp.e)
    x = mp.mpf("0.5") + 2 * mp.sqrt(y / w)
    for _ in range(30):
        step = (logK(x) - y) / dlogK(x)
        x -= step
        if abs(step) < mp.mpf("1e-70"):
            break
    return x


def canonical(y):
    """Return X_0, X_1, X_2, X_3."""
    y = mp.mpf(y)
    w = mp.lambertw(4 * y / mp.e)
    rho = 2 * mp.sqrt(y / w)
    L = mp.log(rho)
    E = C - L / 24
    c1 = -mp.mpf(7) / 5760
    c2 = mp.mpf(31) / 161280

    V1 = -E / L
    V2 = -((L + 1) * V1**2 / 2 - V1 / 24 + c1) / L
    V3 = -((L + 1) * V1 * V2 + V1**3 / 6
           - (V2 - V1**2 / 2) / 24
           - 2 * c1 * V1 + c2) / L

    X0 = mp.mpf("0.5") + rho
    X1 = X0 + V1 / rho
    X2 = X1 + V2 / rho**3
    X3 = X2 + V3 / rho**5
    return X0, X1, X2, X3


def accelerated(y):
    """Return A_0, A_2, A_3, A_4, A_5."""
    y = mp.mpf(y)
    eta = y - C
    omega = mp.lambertw(4 * eta / mp.e)
    tau = 4 * eta / omega
    S = mp.log(tau)

    U2 = 1 / (120 * S)
    U3 = -mp.mpf(31) / (22680 * S)
    U4 = ((1030 * S**2 - 126 * S - 63)
          / (mp.mpf(1814400) * S**3))
    U5 = -((16110 * S**2 - 1023 * S - 341)
           / (mp.mpf(29937600) * S**3))

    series = [
        mp.mpf(1),
        1 + U2 / tau**2,
        1 + U2 / tau**2 + U3 / tau**3,
        1 + U2 / tau**2 + U3 / tau**3 + U4 / tau**4,
        1 + U2 / tau**2 + U3 / tau**3
          + U4 / tau**4 + U5 / tau**5,
    ]
    return tuple(mp.mpf("0.5")
                 + mp.sqrt(mp.mpf(1) / 12 + tau * s)
                 for s in series)


for y in [10, 100, 1000, 10000, mp.mpf("1e6")]:
    exact = inverse_logK(y)
    print("Y =", y)
    print("x =", mp.nstr(exact, 30))
    print("canonical errors:",
          [mp.nstr(x - exact, 12) for x in canonical(y)])
    print("accelerated errors:",
          [mp.nstr(x - exact, 12) for x in accelerated(y)])
\end{lstlisting}

\section{Further deductions and research directions}
\label{sec:future}

\subsection{Late-term growth of the inverse coefficients}

The direct coefficients \(c_n\) inherit factorial growth from Bernoulli
numbers.  The triangular recurrence then forces factorial growth of the
inverse coefficient sequence as well.  A detailed late-term analysis should
identify the Borel singularities and Stokes constants of the inverse directly
from those of \(\log K\).

\begin{conjecture}[Inherited resurgent singularities]
After the natural normalization by the derivative of the dominant phase,
the Borel singularities controlling the late terms of the inverse K
transseries occur at the same action values as those of the centered direct
K expansion.  In particular, the leading action magnitude on the positive
ray is \(2\pi\rho\), and the inverse Stokes constants are obtained from the
direct constants by the transport law in
\cref{thm:exponential-transport}, followed by nonlinear convolution for
higher sectors.
\end{conjecture}

This is strongly suggested by general resurgent inversion principles, but a
complete proof would require a sectorially normalized transseries for
\(\log K\) and careful control of composition in the relevant resurgent
algebra.

\subsection{Optimality of the quadratic coordinate}

Among transformations of the form \(T=r^2+\alpha\), the choice
\(\alpha=-1/12\) is uniquely forced by cancellation of the logarithmic
correction.  A broader optimization problem asks whether a full formal
coordinate
\[
T=r^2-\frac1{12}+\alpha_1r^{-2}+\alpha_2r^{-4}+\cdots
\]
can remove arbitrarily many of the \(d_nT^{-n}\) terms.  Formally, many terms
can be shifted between the coordinate and the residual series, but useful
normal forms should also preserve simplicity, invertibility, and sectorial
analyticity.  The best compromise may be characterized by a canonical
resurgence or Borel-normalization condition.

\subsection{Uniform asymptotics near the Lambert branch point}

The ordinary recurrence divides by \(\Omega+1\), so it is not uniform near
\(\Omega=-1\).  A separate double-scaling theory should combine the
square-root expansion of W near \(-\e^{-1}\) with the lower-order
transseries of \(F\).  Such a theory would describe the merger of inverse
branches and is relevant to global complex continuation.

\subsection{Formalization prospects}

The algebraic core is well suited to proof-assistant formalization:
\begin{itemize}
\item define the \(q\)-adic coefficient ring with a distinguished slow
parameter;
\item formalize the implicit recurrence and uniqueness proof;
\item verify coefficient identities such as \cref{eq:V1,eq:V2,eq:V3};
\item separate these formal results from analytic theorems about Poincare
remainders and monotonicity.
\end{itemize}
The analytic promotion theorem can then be connected to existing formalized
real analysis through explicit residual bounds.

\subsection{Global inverse geometry}

A full account of the inverse K-function should determine:
\begin{enumerate}
\item the critical points and critical values of K;
\item the monodromy relating the \((m,k)\) asymptotic labels;
\item the sectors in which each asymptotic branch is dominant;
\item the Stokes graph for exponentially improved inverse expansions.
\end{enumerate}
The formulas in \cref{sec:complex-branches} provide asymptotic local charts
for this future Riemann-surface description.

\section{Conclusion}

The large inverse of the K-function becomes systematic once its dominant
power--logarithmic phase is inverted exactly.  The unique centering
\(x=r+\tfrac12\) reveals an even asymptotic expansion, and the Lambert
normalizer
\[
\rho=2\sqrt{\frac{\log Z}{W_0(4\log Z/\e)}}
\]
turns inversion into a triangular series in \(\rho^{-2}\).  This yields an
all-orders formula with explicitly recursive coefficients.

The further coordinate
\[
T=x^2-x+\frac16
\]
is not an ad hoc numerical trick.  It is the normal-form transformation that
absorbs the logarithmic correction dictated by the centered expansion.  In
that coordinate, the leading map is
\(\tfrac14T(\log T-1)+C\), its inverse is again exact in terms of W, and the
remaining correction begins two powers later in the relative scale.  The
result is both conceptually cleaner and numerically superior.

More generally, every transseries with dominant part
\(a x^p(\log x-b)\) has an exact W-adapted leading inverse.  Once this scale
is chosen, formal inversion is an implicit-function problem over a slow
coefficient field; residual estimates promote formal identities to analytic
asymptotics; Newton iteration accelerates the expansion; and exponentially
small direct sectors are transported by division by the dominant derivative.
The K-function therefore serves as a particularly transparent model for a
large family of asymptotic inverse problems.

\appendix

\section{Coefficient derivations}
\label{app:coefficients}

\subsection{Centered coefficients}

Starting from \cref{eq:shifted-stirling}, the coefficient of
\(r^{-2n-1}\) in \(\kappa'(r+\tfrac12)\) is
\[
\frac{B_{2n+2}(1/2)}{(2n+1)(2n+2)}.
\]
Integration gives
\[
-\frac{B_{2n+2}(1/2)}{(2n)(2n+1)(2n+2)}r^{-2n},
\]
which is \cref{eq:c-n}.  The identity
\[
B_{2n+2}(1/2)=\left(2^{-2n-1}-1\right)B_{2n+2}
\]
may be used to express all coefficients in terms of ordinary Bernoulli
numbers.

\subsection{Quadratic-normal-form coefficients}

From the proof of \cref{thm:T-normal-form},
\begin{align*}
\kappa(x)-\Psi(T)-C
={}&\frac{T}{4}\log\left(1+\frac{a_0}{T}\right)-\frac{a_0}{4}
+\sum_{m\ge1}c_m(T+a_0)^{-m},
\end{align*}
where \(a_0=1/12\).  The logarithmic part contributes to \(T^{-n}\)
\[
\frac{(-1)^na_0^{n+1}}{4(n+1)}.
\]
Furthermore,
\[
(T+a_0)^{-m}=T^{-m}
\sum_{j\ge0}(-1)^j\binom{m+j-1}{j}a_0^jT^{-j}.
\]
Taking \(j=n-m\) gives
\[
(-1)^{n-m}\binom{n-1}{m-1}a_0^{n-m},
\]
and summing over \(m\le n\) proves \cref{eq:d-n}.

\section{A compact symbolic pseudocode}
\label{app:pseudocode}

The following language-independent pseudocode emphasizes that only truncated
series arithmetic is required.

\begin{lstlisting}[caption={Formal recurrence for the canonical coefficients.}]
INPUT: target order N, symbols L and C
q := formal small parameter
E := C - L/24
v := 0

for n from 1 to N:
    D := (1/2)*(1+v)^2*(L + log(1+v))
         - (1/4)*(1+v)^2 - L/2 + 1/4

    B := E - log(1+v)/24
    for m from 1 to n-1:
        c[m] := -BernoulliPolynomial(2*m+2, 1/2)
                / ((2*m)*(2*m+1)*(2*m+2))
        B := B + c[m]*q^m*(1+v)^(-2*m)

    known := coefficient(truncate(D + q*B, q^(n+1)), q^n)
    V[n] := simplify(-known/L)
    v := v + V[n]*q^n

OUTPUT: V[1], ..., V[N]
CHECK: truncate(D + q*B, q^(N+1)) == 0
\end{lstlisting}

\section{Bibliographic notes}

The K-function definitions and functional identities may be found in
standard special-function references and in the concise summary
\cite{KFunctionWiki}.  Barnes G asymptotics and the Glaisher constant are
recorded in \cite{DLMFBarnesG}; rigorous error bounds and exponential
improvement are developed in \cite{Nemes2014}.  Lambert W and its large
argument expansion are treated in \cite{Corless1996,DLMFW}.  General
transseries background may be found in \cite{vanDerHoeven2006}, while
classical asymptotic error methods are covered by \cite{Olver1997}.  Recent
hyperfactorial refinement literature includes \cite{Xu2020}.

\begin{thebibliography}{99}

\bibitem{Adamchik1998}
V. S. Adamchik,
\newblock PolyGamma functions of negative order,
\newblock \emph{Journal of Computational and Applied Mathematics}
100 (1998), 191--199.

\bibitem{Barnes1900}
E. W. Barnes,
\newblock The theory of the G-function,
\newblock \emph{Quarterly Journal of Pure and Applied Mathematics}
31 (1900), 264--314.

\bibitem{Corless1996}
R. M. Corless, G. H. Gonnet, D. E. G. Hare, D. J. Jeffrey, and D. E. Knuth,
\newblock On the Lambert W function,
\newblock \emph{Advances in Computational Mathematics} 5 (1996), 329--359.

\bibitem{DLMFBarnesG}
NIST Digital Library of Mathematical Functions,
\newblock Section 5.17, Barnes' G-function (Double Gamma Function),
\newblock release 1.2.7, 15 June 2026,
\newblock \url{https://dlmf.nist.gov/5.17}.

\bibitem{DLMFW}
NIST Digital Library of Mathematical Functions,
\newblock Section 4.13, Lambert W-function,
\newblock release 1.2.7, 15 June 2026,
\newblock \url{https://dlmf.nist.gov/4.13}.

\bibitem{KFunctionWiki}
Wikipedia contributors,
\newblock K-function,
\newblock \emph{Wikipedia, The Free Encyclopedia}, accessed 3 September 2026,
\newblock \url{https://en.wikipedia.org/wiki/K-function}.

\bibitem{Nemes2014}
G. Nemes,
\newblock Error bounds and exponential improvement for the asymptotic
expansion of the Barnes G-function,
\newblock \emph{Proceedings of the Royal Society A} 470 (2014), 20140534.

\bibitem{Olver1997}
F. W. J. Olver,
\newblock \emph{Asymptotics and Special Functions},
\newblock A K Peters, Wellesley, 1997 reprint.

\bibitem{ProveItSeries}
V. Reshetnikov,
\newblock Series and transseries research drafts,
\newblock ProveIt repository, 2026,
\newblock \url{https://github.com/VladimirReshetnikov/ProveIt/tree/main/Analysis/FabiusFunction/docs/semi-formalized-research-frontiers/drafts/series-and-transseries}.

\bibitem{vanDerHoeven2006}
J. van der Hoeven,
\newblock \emph{Transseries and Real Differential Algebra},
\newblock Lecture Notes in Mathematics 1888, Springer, Berlin, 2006.

\bibitem{Xu2020}
J. Xu,
\newblock New asymptotic expansions on hyperfactorial functions,
\newblock \emph{Comptes Rendus Mathematique} 358 (2020), no. 9--10,
971--980.

\end{thebibliography}

\end{document}
